\documentclass[spanish,xcolor=dvipsnames]{beamer}
    \usefonttheme{professionalfonts}
    \usetheme{Copenhagen}
    \usecolortheme[named=NavyBlue]{structure}
    \setbeamersize{text margin left=1em,text margin right=1.5em}
    \setbeamercovered{transparent}
    \setbeamertemplate{sections/subsections in toc}[circle]
    %\setbeamertemplate{itemize items}[bullet]
    %\setbeamertemplate{itemize subitem}[square]
    %\setbeamertemplate{enumerate items}[square]
    \setbeamertemplate{navigation symbols}{}
    \setbeamertemplate{headline}{}
	
    \usepackage[spanish]{babel}
    \usepackage[T1]{fontenc}
    \usepackage[utf8]{inputenc}
    
	%Estos paquetes pueden ser útiles,
	%si no están instalados puede quitarlos,
%%	%si los necesita instálelos
    \usepackage{ragged2e}
    \usepackage{lmodern}
    \usepackage[font=scriptsize,labelfont=bf,center]{caption}
    \usepackage{float}
    \usepackage{graphicx}
    \usepackage{subfig}
    \usepackage{multicol}
    \usepackage{amsmath}
    \usepackage{tikz}
    \usetikzlibrary{%
                    automata,%
                    arrows,%
                    shapes,%
                    chains,%
                    matrix,%
                    plotmarks,%
                    %pgfplots.units,%
                    decorations.markings,%
                    decorations.pathmorphing,%
                    backgrounds,%
                    scopes,%
                    calc,%
                    petri,%
                    positioning,%
                    fit%
                    }%
    \usepgflibrary{arrows}
    \usepackage{pgfplots}
    \renewcommand{\shorthandsspanish}{}
    \renewcommand{\spanishtablename}{Tabla}
%    \spanishdecimal{,}

    \makeatletter
        \defbeamertemplate*{footline}{jtheme}
        {
          \leavevmode%
          \hbox{%
          \begin{beamercolorbox}[wd=.25\paperwidth,ht=2.25ex,dp=1ex,center]{author in head/foot}%
            \usebeamerfont{author in head/foot}\insertshortauthor~~\beamer@ifempty{\insertshortinstitute}{}
            {(\insertshortinstitute)}
          \end{beamercolorbox}%
          \begin{beamercolorbox}[wd=.50\paperwidth,ht=2.25ex,dp=1ex,center]{title in head/foot}%
            \usebeamerfont{title in head/foot}\insertshorttitle
          \end{beamercolorbox}%
          \begin{beamercolorbox}[wd=.25\paperwidth,ht=2.25ex,dp=1ex,right]{date in head/foot}%
            \usebeamerfont{date in head/foot}\insertshortdate{}\hspace*{2em}
            \insertframenumber{} / \inserttotalframenumber\hspace*{2ex}
          \end{beamercolorbox}}%
          \vskip0pt%
        }
    \makeatother

\DeclareMathOperator*{\argmin}{arg\ min}
\DeclareMathOperator*{\argmax}{arg\ max}

    %%-------------------Portada-------------------------%%
    \title[Diseño e implementación de un Sistema Web para el Simulador de Eventos Discretos GALATEA]{\sc Sistema Web GALATEA}
    \author[Erik Velásquez ]{ Erik Velasquez}
    \institute[ULA]{
        Facultad de Ingeniería\\Escuela de Ingeniería de Sistemas\\Trabajo de Grado
    }
    \date[Diciembre 2016]{\footnotesize Mérida, 09 de diciembre, 2016}

    %%-------------------Fin-Portada---------------------%%

\begin{document}
	%%--------------------Diapositiva--------------------%%
    \begin{frame}[plain] % [plain] para que no salga el contenido en esta página
	    \begin{center}
	    	\includegraphics[width=1.5cm]{ULA_logo_titulo}
	    \end{center}
	    \titlepage
    \end{frame}
    %%------------------------Fin------------------------%%
	
	%Logo de la esquina superior derecha de cada frame, necesita tres compilaciones para actualizar.
	\addtobeamertemplate{frametitle}{}{%
    \begin{tikzpicture}[remember picture,overlay]
    	\node[anchor=north east,yshift=11pt,xshift=11pt] at (current page.north east) {\includegraphics[height=1.5cm,angle=-20,keepaspectratio]{ula}};
    \end{tikzpicture}}
    
	%%--------------------Diapositiva--------------------%%
    \begin{frame}\justifying
        \frametitle{Agenda}
        \begin{columns}
            \begin{column}{0.0cm}
            \end{column}
            \begin{column}{0.9\textwidth}
                \tableofcontents
            \end{column}
        \end{columns}
    \end{frame}
    %%------------------------Fin------------------------%%

    %%--------------------Diapositiva--------------------%%
    \section{Antecedents}
    \begin{frame}
        \frametitle{Target Detection and Classification in Multispectral Imagery via Sparse Representation}
        %%\framesubtitle{Subtítulo del frame (no obligatorio)}
   
   Se piensa en reportar antecedentes sobre Detección y Clasificación para recordar de donde surge la idea de entrenar un Diccionario (esto de forma muy breve). 
  
    \end{frame}
    %%------------------------Fin------------------------%%

 %%--------------------Diapositiva--------------------%%
    \section{Target Detection}
  
    \begin{frame}
    	\frametitle{Target Detection}
    	\framesubtitle{Sparsity Model}
    
    Sea ${\bf y} \in \mathbb{R}^{n}$  un pixel de observacion multiespectral (pixel de una imagen multiespectral). Entonces\\

	\begin{itemize}
    		\item[\color {black}\bf $\bullet$] ${\bf y}$ se representa como un vector $n-dimensional$ cuyos elementos corresponden a la respuesta a varias bandas espectrales\\
	\end{itemize}

    	\begin{equation}
		\nonumber {\bf y}= {\left[ \begin{array}{cccc} {y}_1 & {y}_2 & \ldots & {y}_n \end{array}  \right]}^T 
	\end{equation}
    	
	\begin{itemize}
    		\item[\color {black}\bf $\bullet$] En nuestro caso, se tienen $4$ bandas espectrales $(n=4)$ 
	\end{itemize}
  
    \end{frame}
    
    %%------------------------Fin------------------------%%

 %%--------------------Diapositiva--------------------%%
    %%\section{Dictionary learning for classification}
    \begin{frame}
   	\frametitle{Sparsity Model}
    	\framesubtitle{Background pixel}

 Si ${\bf y}$ es un \textcolor{black} {\bf pixel de fondo}, éste puede ser representado aproximadamente como una combinacion lineal de átomos de fondo ${\{{\bf d}_i\}}_{i=1,2,...,K_b}$  

	\begin{equation}
	\nonumber  {\bf y} \approx {\alpha}_1{\bf d}_1^b+{\alpha}_2{\bf d}_2^b+...+{\alpha}_{ K_b} {\bf d}_{K_b}^b
	\end{equation}

	\begin{equation}
	\nonumber  {\bf y} \approx \underbrace{ \left[ \begin {array}{cccc} 
{\bf d}_1^b & {\bf d}_2^b & \ldots & {\bf d}_{K_b}^b \end {array} \right]}_{{\bf D}_b}
 {\underbrace{ \left[ \begin{array}{cccc} {\alpha}_1 & {\alpha}_2 & \ldots & {\alpha}_{K_b} \end{array} \right]} _{{\bf x}_b}}^T
	\end{equation}

	\begin{equation}
	\nonumber  {\bf y} = {\bf D}_b  {\bf x}_b
    	\end{equation}

Donde:
	\begin{itemize}	
		\item[\color {black}\bf $\bullet$] ${\bf D}_b \in \mathbb{R}^{n \times {K_b}}$: Es el diccionario de fondo.\\
		\item[\color {black}\bf $\bullet$] ${\bf x}_b \in \mathbb{R}^{K_b \times 1}$: Es el vector de coeficientes de fondo.\\
	\end{itemize}
 
    \end{frame}
    %%------------------------Fin------------------------%%


%%--------------------Diapositiva--------------------%%
   %% \section{Dictionary learning for classification}
    \begin{frame}
    	\frametitle{Sparsity Model}
    	\framesubtitle{Target pixel}

 Si ${\bf y}$ es un \textcolor{black} {\bf pixel objetivo}, éste puede ser representado aproximadamente como una combinacion lineal de átomos de objetivo ${\{{\bf d}_i\}}_{i=1,2,...,K_t}$  

	\begin{equation}
		\nonumber  {\bf y} \approx {\beta}_1{\bf d}_1^t+{\beta}_2{\bf d}_2^t+...+{\beta}_{K_t}{\bf d}_{K_t}^t
	\end{equation}

	\begin{equation}
		\nonumber  {\bf y} \approx \underbrace{ \left[ \begin{array}{cccc}
 {\bf d}_1^t & {\bf d}_2^t & \ldots & {\bf d}_{K_t}^t \end{array} \right]}_{{\bf D}_t} 
{\underbrace{ \left[ \begin{array}{cccc} {\beta}_1 & {\beta}_2 & \ldots & {\beta}_{K_t} \end{array} \right]}_{{\bf x}_t}}^T
	\end{equation}

	\begin{equation}
		\nonumber  {\bf y} = {\bf D}_t {\bf x}_t
    	\end{equation}

Donde:
	\begin{itemize}	
		\item[\color {black}\bf $\bullet$] ${\bf D}_t \in \mathbb{R}^{n \times {K_t}}$: Es el diccionario de objetivo.\\
		\item[\color {black}\bf $\bullet$] ${\bf x}_t \in \mathbb{R}^{K_t \times 1}$: Es el vector de coeficientes de objetivo.\\
	\end{itemize}

\end{frame}
%%--------------------Fin Diapositiva--------------------%%

%%--------------------Diapositiva--------------------%%
%%\section{Label consistent K-SVD}

\begin{frame}
	\frametitle{Sparsity Model}
	\framesubtitle{Global Pixel}

Para modelar cualquier pixel de observación ({\bf pixel global}), se combinan los diccionarios ${\bf D}_b$ y ${\bf D}_t$ formando así un {\bf diccionario global} ${\bf D}$. Por tanto, un {\bf pixel global} de prueba se puede escribir como una combinación lineal de unos pocos {\bf átomos globales} ${\{{\bf d}_i\}}_{i=1,2,...,K}$. Así

	\begin {equation}
		\nonumber	{\bf y}=
		\underbrace{\left[
		\begin {array}{cc}
		{\bf D}_b & {\bf D}_t
		\end {array} 
		\right]}_{\bf D}
		\underbrace{\left[
		\begin {array}{c}
		{\bf x}_b \\ {\bf x}_t
		\end {array} 
		\right]}_{\bf x}
	\end{equation}	

	\begin {equation}
		\nonumber {\bf y}={\bf D}{\bf x}
	\end {equation}

Donde:
	\begin{itemize}	
		\item[\color {black}\bf $\bullet$] ${\bf D} \in \mathbb{R}^{n \times {K}}$: Es el diccionario global (con $K=K_b+K_t$).\\
		\item[\color {black}\bf $\bullet$] ${{\bf x}} \in \mathbb{R}^{K \times 1}$: Es el vector de coeficientes global.\\
	\end{itemize}

\end{frame}
%%--------------------FIN Diapositiva--------------------%%

%%--------------------Diapositiva--------------------%%
%%\subsection{LC-KSVD1}
\begin{frame}
	\frametitle{Reconstruction Task}
%%\framesubtitle{LC-KSVD1}

Asumiendo que se conoce el diccionario global ${\bf D}$, ${\bf x}$ se puede hallar resolviendo el problema de minimización (ampliamente conocido)

	\begin {equation}
		\nonumber {\bf x}=\argmin_{\bf x}\|{\bf y}-{\bf Dx}\|_2^2 \; s.t. \; \|{\bf x}\|_0 \leq T
	\end {equation}

Una vez conocido (reconstruido) el vector de coeficientes ${\bf x}$, se pueden determinar los respectivos ${\bf x}_b$ y ${\bf x}_t$

	\begin {equation}
		\nonumber	{\bf x}=
		\left[
		\begin {array}{c}
		{\bf x}_b \\ {\bf x}_t
		\end {array} 
		\right]
	\end {equation}

\end{frame}	

%%--------------------Fin Diapositiva--------------------%%

%%--------------------Diapositiva--------------------%%

\begin{frame}
	\frametitle{Detection Task}
	%%\framesubtitle{LC-KSVD1}

Con los correspondientes vectores de coeficientes ${\bf x}_b$ y ${\bf x}_t$, se pueden calcular los respectivos errores de reconstruccón

	\begin {equation}
		\nonumber {\bf r}_b ( {\bf y} )=\| {\bf y} - {\bf D}_b  {\bf x}_b \|_2
	\end{equation}
	\begin {equation}
		\nonumber {\bf r}_t ( {\bf y} )=\| {\bf y} - {\bf D}_t  {\bf x}_t \|_2
	\end{equation}

Luego, el detector podría plantearse de la sig. manera

	\begin {equation}
		\nonumber {\bf TD} ( {\bf y} )={\bf r}_b ( {\bf y} ) - {\bf r}_t ( {\bf y} )
	\end{equation}

Donde:
	\begin{itemize}	
		\item[\color {black}\bf $\bullet$] Si ${\bf TD} ( {\bf y} ) > \delta$: Entonces ${\bf x}$ es determinado como un pixel objetivo.\\
		\item[\color {black}\bf $\bullet$] Si ${\bf TD} ( {\bf y} ) < \delta$: Entonces ${\bf x}$ es determinado como un pixel de fondo.\\
	\end{itemize}	
	
\end{frame}
%%--------------------Fin Diapositiva--------------------%%

%%--------------------Diapositiva--------------------%%

\begin{frame}
	\frametitle{Observations}
	\framesubtitle{First}

Para {\bf D} se trabajó inicialmente con átomos tomados directamente de muestras sobre la imagen.	

	\begin {itemize}
		\item [\color {black} \bf $\bullet$] ${\bf D}_b:$ átomos conformados de pixeles de fondo, directamente.
		\item [\color {black} \bf $\bullet$] ${\bf D}_t:$ átomos conformados de pixeles de objetivo, directamente.
	\end {itemize}

Lo cual, probablemente, conduce a una reconstrucción relativamente pobre (menos discriminativa).

	\begin {itemize}
		\item [\color {red} \bf $\blacktriangleright$] \textcolor{red} {Por esta razón, se pensó en el entrenamiento de un diccionario. Apareciendo el estudio del algoritmo K-SVD como mejor opción.}
	\item [\color {red} \bf $\blacktriangleright$] \textcolor{red} {Con lo que se espera una reconstrucción más discriminativa, que potencie el rendimiento del detector.}
	\end {itemize}

\end{frame}	

%%--------------------Fin Diapositiva--------------------%%



 %%--------------------Diapositiva--------------------%%
 
 \begin{frame}
 	\frametitle{Observations}
 	\framesubtitle{Second}
 
Con sólo entrenar un diccionario ${\bf D}$ que se adapte a la data (mediante K-SVD) se procedería con las tareas de reconstrucción y luego detección. Con ello:

	\begin {itemize}
		\item [\color {black} \bf $\bullet$] Se espera un mejor rendimiento en la detección debido a una reconstrucción de más fiable.
		\item [\color {black} \bf $\bullet$] No obstante, los procesos de {\bf entrenamiento} y {\bf reconstrucción} se trabajan aisladamente.
	\end {itemize}

	Lo anterior, conduce a pensar en algo más ambisioso.

	\begin {itemize}
		\item [\color {red} \bf $\blacktriangleright$] \textcolor {red} {Durante el mismo entrenamiento, se trabajen ambas operaciones (entrenamiento/detección) de manera simultánea.}
	\end {itemize}
 	
 \end{frame}
 
 %%------------------------Fin------------------------%%


%%--------------------Diapositiva--------------------%%
\section {Classification}
\begin {frame}
	\frametitle {Classification}

Suponiendo que se presentan \textit {\textbf{m}} posibles clases.
El diccionario global estaría determinado por la combinación de \textit {\textbf{m}} subdicionarios. 
(uno para clada clase)

	\begin{equation}
		\nonumber {\bf D}=\left[ \begin{array}{cccc} {\bf D}_1 & {\bf D}_2 & \ldots & {\bf D}_m \end{array} \right]
	\end{equation}

Para lo cual, resolviendo la tarea de reconstrucción, se obtendría un vector de coeficientes global ${\bf x}$ formado por la combinación de \textit {\textbf{m}} subvectores de coeficientes (uno por clase).

	\begin{equation}
		\nonumber {\bf x}=\left[ \begin{array}{c} {\bf x}_1 \\ {\bf x}_2 \\  \vdots \\ {\bf x}_m \end{array} \right]
	\end{equation}

\end{frame}
 %%--------------------Fin--------------------%%


   %%--------------------Diapositiva--------------------%%
\begin {frame}
	\frametitle{Classification}

Con los \textit {\textbf{m}} subdicionrios y vectores de coeficientes, se pueden calcular \textit {\textbf{m}} errores de reconstrucción (uno para cada clase).

	\begin{equation}
		\nonumber {\bf r}_i ({\bf y}) = \| {\bf y} - {\bf D}_i{\bf x}_i \|_2 \  \; \; \; \; \; \; {}_{i=1,2,\ldots, m}
	\end {equation}

Obteniéndose
	\begin{equation}
		\nonumber {\bf r} (i) = \left[ \begin{array}{c} {\bf r}_1 ({\bf y}) \\ {\bf r}_2 ({\bf y}) \\ \vdots \\ {\bf r}_m ({\bf y}) \end{array} \right] 
	\end {equation}

Luego, la clase queda determinada por

	\begin{equation}
		\nonumber {\bf C} \textnormal{lass} ({\bf y}) = i = \argmin_i \ {{\bf r} (i)}  \; \; \; \; \; \;  {}_{i=1,2,\ldots, m}
	\end{equation}
		
\end {frame}

%% ----------------------New Frame--------------------%%
\section {SVD}
\begin{frame}
	\frametitle{Singular Value Descomposition}
	\framesubtitle{Definition}
		
Cualquier matriz rectangular ${\bf A} \in \mathbb{R}^{m \times n}$ se puede factorizar de la forma

	\begin{equation}
		\nonumber {\bf A}={\bf U} {\bf \Sigma} {\bf V}^T
	\end{equation}

Donde:

	\begin{itemize}
		\item[\color{black}  \bf $\bullet$] ${\bf U} \in \mathbb{R}^{m \times m}:$ Matriz cuyas columnas son los autovectores de ${\bf A}{\bf A}^T$ 
		\item[\color{black}  \bf $\bullet$] ${\bf V} \in \mathbb{R}^{n \times n}:$ Matriz cuyas columnas son los autovectores de ${\bf A}^T{\bf A}$
		\item[\color{black}  \bf $\bullet$] ${\bf \Sigma} \in \mathbb{R}^{m \times n}:$ Matriz diagonal (pero rectangular) que contiene la raiz cuadrada de los autovalores no-negativos de ${\bf V}$. \newline Ellos son los valores singulares de ${\bf A}$
	\end{itemize}
\end{frame}	
%%-----------------------End Frame--------------------%%

%%----------------------New Frame--------------------%%
\begin{frame}
	\frametitle{SVD}
	\framesubtitle{Particularity}
	
Segun su definición, SVD es una multiplicación de matrices.

	\begin{itemize}
		\item [\color {black} \bf $\bullet$] Se puede manejar esta multiplicacion como $columns \ times \ rows$
	\end{itemize}

	\begin{equation}
		\nonumber {\bf A}={\bf U} {\bf \Sigma} {\bf V}^T={\bf u}_1{\bf \sigma}_1{\bf v}_1^T+{\bf u}_2{\bf \sigma}_2{\bf v}_2^T+ \ldots + {\bf u}_r{\bf \sigma}_r{\bf v}_r^T
	\end{equation}

	\begin{itemize}
		\item [\color{red} \bf $\blacktriangleright$] \textcolor{red} {Cualquier matriz ${\bf A}$ (rectangular) se puede \textbf{descomponer} como la suma de $r$ submatrices de $rango$$-$$1$}.
		\item [\color{red} \bf $\blacktriangleright$] \textcolor{red} {Cada matriz está formada por la multiplicacion entre la $i$$-$$th$ columna de ${\bf U}$ y la $i$$-$$th$ fila de ${\bf V}^T$, además de estar ponderadas/escaladas/pesadas por el correspondiente autovalor (\textbf{valor singular})}.
	\end{itemize}

\end{frame}
%%------------------------End Frame-----------------%%
 
%%------------------------New Frame----------------%%
\section{K-SVD}
\begin{frame}
	\frametitle{Algorithm K-SVD}
	\framesubtitle{Definition And Operability}
Dadas un conjunto de $N$ señales de entrenamineto ${\bf Y}$, el algoritmo busca el diccionario ${\bf D}$ que conduce a la mejor representacion ${\bf X}$ de cada uno de los $N$ miembros en ese conjunto, bajo una estricta $sparsity \ constraint$.

	\begin{equation} \label{eq:main problem}
		<{\bf D},{\bf X}>=\argmin_{\bf D,\bf X}\{\|{\bf Y}-{\bf DX}\|\}_F^2\  s.t.\ \forall i, \|{\bf x}_i\|_0\leq T    	\end{equation}

Para tal fin, trabaja en base a dos etapas

	\begin{enumerate}
		\item $Sparse \ Coding \ Stage$.
		\item $Process \ of \ Updating \ the \ dictionary$.
	\end{enumerate}

\end{frame}
%%--------------------------End Frame----------------------------%%

%%--------------------------New Frame----------------------------%%
\begin {frame}
	\frametitle{First Stage: Sparse Coding Stage}

	\begin{itemize}
		\item[\color {black} \bf $\bullet$] Se asume un ${\bf D}$ fijo y se considera el problema de optimización (\ref{eq:main problem}).
		\item[\color {black} \bf $\bullet$] Se halla la matriz de de representación poco densa ${\bf X}$ ($Sparse$  $Coding$)  %\newline
	\end{itemize}
Esto se logra, resolviendo $N$ problemas

	\begin{equation}
		\nonumber {\bf x}_i=\argmin_{{\bf x}_i}\|{\bf y}-{\bf Dx}_i\|_2^2 \; \; s.t. \; \|{\bf x}\|_0 \leq T
	\end{equation}
		
Donde la $i$$-$$th$ columna en ${\bf X}$, corresponde al vector de codificación ${\bf x}_i$ de la $i$$-$$th$ señal de entrenamineto ${\bf y}_i$ en ${\bf Y}$ para el ${\bf D}$ fijo. \newline 

\end {frame} 
%%--------------------------End Frame----------------------------%%

%%--------------------------New Frame----------------------------%%
\begin {frame}
	\frametitle{Second Stage: Process of Updating the dictionary}

	\begin{itemize}
		\item[\color {black} \bf $\bullet$] Se asume que tanto ${\bf D}$ como ${\bf X}$ están fijos y se considera el problema de optimización (\ref{eq:main problem}).
	\end{itemize} 

 Recordando la multiplicación matricial vista como $columns \ times \  rows$. Entonces

	\begin{equation}
		\nonumber {\bf DX} = \sum_{j=1}^K {\bf d}_j {\bf x}_T^j
	\end{equation}
	
 Con lo cual, el término $\| {\bf Y} - {\bf DX} \|_F^2$ en (\ref{eq:main problem}) se puede reescribir como
	
	\begin{equation} \label{eq:auxiliar}
		\| {\bf Y} - {\bf DX} \|_F^2 = \left\| {\bf Y} - \sum_{j=1}^K {\bf d}_j {\bf x}_T^j  \right\|_F^2
	\end{equation}
	
\end {frame}

%%--------------------------End Frame----------------------------%%

%%--------------------------New Frame----------------------------%%
\begin {frame}
	\frametitle{Second Stage: Process of Updating the dictionary}

	\begin{itemize}
		\item[\color {black} \bf $\bullet$] Para proponer la actualización de ${\bf D}$, se pone en cuestion un sólo átomo ${\bf d}_k$ y la fila de coeficientes ${\bf x}_T^k$ correspondientes
	\end{itemize} 

Ello se logra reescribiendo (\ref{eq:auxiliar}) de la sig. manera

	\begin{equation} 
		\nonumber \left \| {\bf Y} - {\bf DX} \right \|_F^2 = \left\| \left( {\bf Y} - \sum_{j \ne k} {\bf d}_j {\bf x}_T^j \right) - {\bf d}_k {\bf x}_T^k \right\|_F^2  \end{equation}
	\begin{equation} \label{eq:dif}
 		= \left \| {\bf E}_k - {\bf d}_k {\bf x}_T^k \right \|_F^2.
	\end{equation}

	\begin {itemize}
		\item [\color {red} \bf $\blacktriangleright$] \textcolor{red} {Se desea minimizar la diferencia en (\ref{eq:dif}) sabiendo que ${\bf d}_k{\bf x}_T^k$ es una matriz de $rango$$-$$1$. ¿Qué se sugiere?}
	\end{itemize}
\end{frame}
%%--------------------------End Frame----------------------------%%

%%--------------------------New Frame----------------------------%%
\begin{frame}
	\frametitle{Second Stage: Process of Updating the dictionary}
	
	\begin {itemize}
		\item [\color {red} \bf $\blacktriangleright$] \textcolor{red} {Aplicar SVD a ${\bf E}_k$ para determinar la matriz de $rango$$-$$1$ más semejante a ella.}
		\item [\color {red} \bf $\blacktriangleright$] \textcolor{red} {Esto se realiza $K$ veces, a favor de actualizar cada uno de los $K$ átomos en ${\bf D}$.} 
		\item [\color {red} \bf $\blacktriangleright$] \textcolor{red} {Lo anterior justifica el nombre del algoritmo, pues se aplica "SVD" $K$ veces (K-SVD)}
	\end{itemize}


\end{frame}
%%--------------------------End Frame----------------------------%%


%%--------------------------New Frame----------------------------%%
\begin{frame}
	\frametitle{Second Stage: Process of Updating the dictionary}
	
Muy importante, realmente (\ref{eq:dif}) se redefine antes de aplicar SVD a favor de mantener una representación poco densa. Así 

	\begin{equation}
		\nonumber \left \| {\bf E}_k {\bf \Omega}_k - {\bf d}_k {\bf x}_T^k {\bf \Omega}_k \right \|_F^2
		= \left \| {\bf E}_k^R - {\bf d}_k {\bf x}_R^k \right \|_F^2.		
	\end{equation} 

Donde:
	\begin {itemize}
		\item [\color {black} \bf $\bullet$]  ${\bf \Omega}_k$ es la matriz que permite preservar el soporte del vector fila ${\bf x}_T^k$
		\item [\color {black} \bf $\bullet$]  ${\bf x}_R^k$ son los coeficientes no nulos de ${\bf x}_T^k$
		\item [\color {black} \bf $\bullet$] $ {\bf E}_k^R$ son las columnas de errores que corresponden a los ejemplos que usan el átomo ${\bf d}_k$
	\end{itemize}

Luego, Se aplica SVD a ${\bf E}_k^R$, con lo cual	
	\begin{equation}
		\nonumber {\bf d}_k {\bf x}_R^k =  {\bf u}_i {\sigma}_i {\bf v}_i^T  
	\end {equation}

\end{frame}
%%--------------------------End Frame----------------------------%%

    %%--------------------Diapositiva--------------------%%
    \section{Bibliography}
    \begin{frame}
        \frametitle{Bibliography}
        \begin{thebibliography}{9}\small
            \setbeamertemplate{bibliography item}[text]
            \bibitem{A}Y. Chen, N. M. Nasrabadi, and T. D. Tran, “Sparse representation for target detection in hyperspectral imagery,” IEEE J. Sel. Topics Signal Process, vol. 5, no. 3, pp. 629-640, June 2011.
            \bibitem{B}Y. Chen, N. M. Nasrabadi, y T. D. Tran, “Hyperspectral Image Classifications using Dictionary-based sparse representation,” IEEE Trans. Geosci. Remote Sens., vol. 49, no. 10, pp. 3973-3985., Oct. 2011.
            \bibitem{C} G. Strang, \textit{Linear Algebra and Its Aplications.} Fourth Edition. 
	 \bibitem{D} M. Aharon, M. Elad, y A. Bruckstein, “K-SVD: An algorithm for designing overcomplete dictionaries for sparse representation,” IEEE Transactions on Signal Processing, vol. 54, no. 11, pp. 4311-4322, Nov. 2006.
        \end{thebibliography}
    \end{frame}
    %%------------------------Fin------------------------%%
\end{document}